% =====================================================================
% Plantilla de examen con banco de problemas y soluciones conmutables.
%
% - Cada problema es un \item con \label{item:<nombre>} y su solución en el
%   entorno solution, justo debajo.
% - Los problemas que no van en el examen se dejan en el banco, dentro de
%   \begin{comment} ... \end{comment}.
% - Para imprimir el examen: deja \excludecomment{solution} (así viene).
%   Para ver las soluciones: comenta esa línea.
% - Para mostrar también lo que va dentro de \phantom (tablas que el
%   estudiante completa), descomente la redefinición de \phantom.
% =====================================================================
\documentclass[12pt, a4paper]{article}

\usepackage[utf8]{inputenc}
\usepackage[T1]{fontenc}
\usepackage[spanish,es-nodecimaldot]{babel}
\usepackage{amsmath, amssymb}
\usepackage{graphicx}
\usepackage[margin=2cm]{geometry}
\usepackage{siunitx}
\usepackage{cancel}
\usepackage{enumitem}
\usepackage{sectsty}
\usepackage{tikz}
\usetikzlibrary{babel, patterns}

% soluciones y banco de problemas
\usepackage{comment}
\specialcomment{solution}
  {\begingroup \textbf{Solución: }}{\medskip \endgroup}
\excludecomment{solution}   % comente esta línea para ver las soluciones
\excludecomment{comment}    % banco de problemas que no van en el examen
% \renewcommand{\phantom}[1]{#1}  % descomente junto con la línea de solution
                                  % para mostrar lo que va dentro de \phantom

\pagestyle{empty}
\sectionfont{\centering\large\bfseries}

\begin{document}

\vspace*{-2cm}
Nombre:$\underline{\hspace{8.5cm}}$ Identificación: \hrulefill

\begin{enumerate}

% ===== Problemas del examen (en este orden) =====
\item\label{item:ejemplo} Una piedra se suelta desde el reposo a una altura $h=\qty{20}{\m}$. Despreciando la fricción con el aire, calcule la rapidez con que llega al suelo. Tome $g=\qty{9.8}{\m/\s^2}$.

\vspace{4cm}

\begin{solution}
Por conservación de la energía mecánica,
\begin{align}
    \cancel{m}gh = \frac{1}{2}\cancel{m}v^2
    \;\Longrightarrow\;
    v = \sqrt{2gh} = \sqrt{2\times 9.8\times 20}\,\unit{\m/\s} \approx \qty{19.8}{\m/\s}\,.
\end{align}
\end{solution}

\newpage
% ===== Banco de problemas no seleccionados =====
\begin{comment}
\item\label{item:banco-ejemplo} Problema del banco: no aparece en el examen.
\begin{solution}
Solución del problema del banco.
\end{solution}
\end{comment}

\end{enumerate}

\end{document}
